\needspace{4\baselineskip}
\color{uniblue}{\subsection{\raggedright{Управление выключателем (УВ)}}}
\color{black}

\begin{enumerate}[label=\arabic{section}.\arabic{subsection}.\arabic*, labelsep=4pt, leftmargin=0pt, itemindent=57pt]

%%%%%%%%%%%%%%%%%%%%%%%%%%%%%%%%%%%%%%%%%%%%%%%%%%%%%%%%%%%%%%%%%%%%%%  ОБЩИЕ СВЕДЕНИЯ %%%%%%%%%%%%%%%%%%%%%%%%%%%%%%%%%%%%%%%%%%%%%%%%%%

\item
Функциональный блок <<УВ>> предназначен для передачи команды на включение выключателя от контроллера присоединения, в случае, когда контроллер присоединения является внешним устройством. В функциональном блоке контролируются условия разрешения для выполнения команды включения выключателя.

\item
Ввод функционального блока <<УВ>> в работу на этапе параметрирования устройства осуществляется программным переключателем SGF1 <<Ввод\_функции>>. Введенное состояние функционального блока соответствует активному состоянию сигнала <<УВ: Ввод>>.
\item
Функция формирует сигнал <<УВ / УВ: Включить>> при активации сигнала <<Опер. включить В>> (команда включения от внешнего КП) и при отсутствии сигнала блокировки включения выключателя от функции контроля выключателя <<КСВ>>. Сигнал оперативного отключения от внешнего КП не требует проверки каких-либо условий и передается напрямую в схему управления выключателем. 
\item 
Функциональный блок <<УВ>> может быть выведен из работы оперативно путем активации сигнала <<ОВ УВ>>, о чем сигнализирует активное состояние сигнала <<УВ: Оперативный вывод>>.
\item
Логика управления режимами работы функционального блока <<УВ>> выполнена в соответствии с рисунком \ref{uv:img_gen}. В таблице \ref{uv:tbl_gen} приведены параметры, необходимые для настройки функционального блока <<УВ>>. 

\vspace{3mm}
\begin{figure}[!h]
\centering
\includegraphics[width=1\textwidth,height=1\textheight,keepaspectratio]{\fbpath/51. УВ/_latex/img/uv.pdf}
\captionof{figure}{Функциональная схема <<УВ>>}
\label{uv:img_gen}
\end{figure}

\small
\begin{longtable}{|>{\centering\arraybackslash}m{5.3cm}|>{\centering\arraybackslash}m{3.3cm}|>{\centering\arraybackslash}m{4.2cm}|>{\centering\arraybackslash}m{1.8cm}|>{\centering\arraybackslash}m{1cm}|}
\caption{Параметры для настройки функционального блока <<УВ>>\hfill\vspace{-0.5\baselineskip}}\label{uv:tbl_gen}\\ 
\hline
\rowcolor{gray!30}
Параметр (Параметр на ИЧМ) & Условное обозначение на схеме & Значение/ Диапазон & Единица измерения & Шаг \\ 
\hline
\endfirsthead
%\caption{\emph{Продолжение\hfill\vspace{-0.5\baselineskip}}} \\ % Отступ обозначения таблицы от таблицы и выравнивание надписи слева
\caption*{\hspace{3pt}\emph{Продолжение таблицы \ref{uv:tbl_gen}\hfill\vspace{-0.5\baselineskip}}} \\ % сделано по ГОСТ 2.105 п.6.8.7
\hline
\rowcolor{gray!30}
Параметр (Параметр на ИЧМ) & Условное обозначение на схеме & Значение/ Диапазон & Единица измерения & Шаг \\ 
%\hline
\endhead
%\hline
\endfoot
%\hline
\endlastfoot
%==+t1*CBCSWI1|T_HVBCTRL
%===t1*CBCSWI1
\centering Ввод функции в работу (Ввод\_функции) & \centering SGF1 & \centering 0 = Не предусмотрено\\1 = Предусмотрено & \centering -- & \centering \arraybackslash -- \\
\hline
%===t1
\end{longtable}
\normalsize

\end{enumerate}
\FloatBarrier % Все float'ы ДО этой строки будут размещены здесь.